\section{Introduction}
Nuclear energy is an important part of the US energy mix and 
will be needed to help meet climate goals. The current 
fleet of US commercial reactors are all \glspl{LWR}, 
which require uranium enriched to less than 5\% $^{235}$U
based on \gls{NRC} regulations \cite{noauthor_10_2006}.
But as energy needs and economics change, many companies 
have developed advanced reactors that will require 
uranium enriched to between 5-20\% $^{235}$U, or \gls{HALEU}. 
Using \gls{HALEU} in reactors, including \glspl{LWR}, 
helps in achieving higher burnups and longer cycle 
times \cite{carlson_implications_2022}
However, there is no commercial way to produce \gls{HALEU} 
for the US right now. 

There are two primary methods for producing \gls{HALEU}. The first 
option is to enrich uranium to the required level. Developing 
infrastructure to enable this \gls{HALEU} production method 
is underway 
\cite{nuclear_energy_institute_establishing_2022,us_nuclear_regulatory_commission_centrus_2021},
but will take time. 
The other \gls{HALEU} production option is to downblend \gls{HEU} to 
the required \gls{HALEU}
level. The \gls{DOE} is considering multiple sources of \gls{HEU} 
to create \gls{HALEU} \cite{nuclear_energy_institute_establishing_2022}. The 
first is spent fuel from \gls{EBR}, which can produce a total of 10 MTU 
of \gls{HALEU} at 19.75\% enrichment \cite{nuclear_energy_institute_establishing_2022}. 
The next source is the stored \gls{HEU} solution at \gls{SRS} from 
spent research reactor fuel. Estimates of the amount of \gls{HALEU}
that can be produced from this stockpile vary from 4 MTU 
\cite{nuclear_energy_institute_establishing_2022} to 20 MTU \cite{regalbuto_addressing_2020}.
These quantities are much less than potential \gls{HALEU} demand 
to support a fleet of \gls{HALEU}-fueled reactors 
\cite{dixon_estimated_2022,bachmann_enrichment_2021}, but the \gls{HEU} 
sources can be used to provide fuel for demonstration 
projects or initial core loadings. 

In addition to the limited quantities of \gls{HEU} available for 
downblending, the \gls{HALEU} produced from downblended \gls{HEU}
has uranium impurities that are not typically included in reactor 
physics analysis \cite{vaden_isotopic_2018,nelson_foreign_2010}. 
Other uranium isotopes present in this \gls{HALEU} includes
$^{234}$U and $^{236}$U, which can affect the neutronics of a reactor. 
\gls{HALEU} produced through enriching natural uranium for use 
in the \gls{MIT} Reactor 
contained measurable concentrations of $^{233}$U, $^{234}$U, $^{235}$U,
$^{236}$U, and $^{238}$U \cite{mascolino_impact_2022}.
Variations in the concentrations of these isotopes affected 
the \keff of the reactor by 280 pcm \cite{mascolino_impact_2022}. 
This is a relatively small change in the \keff of a reactor, but the 
results show that the concentrations of different uranium isotopes 
does affect performance of a reactor. 
Other analysis analyzed the effects of the uranium composition 
for a potential research reactor at \gls{NIST}, comparing the 
performance of nominally enriched \gls{HALEU} with 
\gls{HALEU} from downblended \gls{HEU} from the Y-12 National 
Security Complex \cite{celikten_effects_2021}. This analysis 
showed a 0.43\% drop in reactivity, which leads to an 8\% 
reduction in cycle length for the reactor \cite{celikten_effects_2021}.
This analysis further shows how the impurities from downblended 
\gls{HEU} can affect reactor performance. 

These previous investigations were both performed based on research 
reactor designs, which are typically much smaller in power and 
size than commercial reactors. Therefore, if downblended 
\gls{HEU} is to be used in commercial reactors, analysis is needed on 
how these impurities may impact performance of a commercial 
reactor design. The purpose of this work is to investigate how the 
impurities in \gls{HALEU} produced from downblending \gls{HEU} can 
affect the neutronics and operation of select advanced reactors 
intended for commercial operation. This analysis will provide 
insight into how the impurities in the \gls{HALEU} fuel affects 
the reactors, and if the impurities are sufficient to potentially 
prevent safe operation of the reactor. 
